\documentclass[11pt,a4paper]{article}
\usepackage[T1]{fontenc}
\usepackage{amsmath,amssymb} % math packages
\usepackage[colorlinks=true]{hyperref}
\makeatletter
\renewcommand\@seccntformat[1]{\csname the#1\endcsname.\quad}
\def\my@helper#1{\textbf{#1}}
\makeatother
%% A full-line comment with $math$ & {braces} [brackets] (parens)
\newcommand{\vect}[1]{\mathbf{#1}}
\newenvironment{note}{\begin{quote}\itshape}{\end{quote}}

\title{Syntax Coloring Sample}
\author{A. Author\thanks{Supported by \emph{nothing}.}}
\date{\today}

\begin{document}
\maketitle

\section{Introduction}\label{sec:intro}
Plain text with an escaped percent \% and an escaped dollar \$ and \& too.
Inline math $a^2 + b^2 = c^2$, display math \[ \int_0^1 f(x)\,dx \] and \( x \) here.
$$ \sum_{k=1}^{n} k = \frac{n(n+1)}{2} $$
A forced break\\ and \\[2pt] a spaced one. Accents: \'e \"o \^{a} \~n \c{c}.
Ties~and dashes -- --- and quotes ``quoted'' text.

\subsection*{Citations, footnotes and indices}
See \cite{knuth84}, \cite[p.~3]{lamport94}, \citep{a,b} and \citet[see][ch.~2]{c}.
Also \autocite[12]{d}, \Cite{e}, \footcite{f} and \parencite{g}.
A footnote\footnote{With {nested {braces}} and $x$ inside.} and another\footnote[3]{Numbered.}.
An index entry\index{coloring!syntax|textbf} and \index[names]{Koch, Richard}.
A Spanish one\esindex{coloreado} and \indexentry{not this}.

\begin{itemize}
  \item First item with \textit{italic} and \texttt{\textbackslash verb}.
  \item[$\star$] Second item & an ampersand.
\end{itemize}

\begin{tabular}{|l|c|r|}
  \hline
  a & b & c \\ \hline
  1 & 2 & 3 \\
\end{tabular}

\begin{align}
  f(x) &= (x+1)^2 \label{eq:f} \\
       &= x^2 + 2x + 1 \tag{$*$}
\end{align}

Verbatim: \verb|%$&{}| and \verb+\foo{bar}+ and \verb*|a b|.
\begin{verbatim}
\notacommand % not a comment either { [ (
\end{verbatim}

Spelling: \ref{sec:intro}, \pageref{eq:f}, \label{x:y}, \emph{checked words}, \textbf{bold words}.
\setcounter{page}{3}\addtocounter{section}{-1}\hspace*{1em}\vspace{2ex plus 1fil}
\includegraphics[width=0.5\linewidth,angle=90]{figures/plot.pdf}
\usepackage[utf8]{inputenc}\input{chapter1}\include{chapter2}
\begin{figure}[htbp]\centering\caption[Short]{Long caption}\end{figure}

Odd commands: \@ \_ \# \{ \} \, \; \! \  \\* \@@ \a@b \x_y:z \foo.bar
Command at line end \LaTeX
Trailing backslash at end of line \
\end{document}
